% ============================================================
\section{Robustez}\label{sec:robustez}
% ============================================================

Os resultados só valem se não dependerem de uma escolha particular de construção. Refizemos as
redes mudando uma escolha de cada vez e repetimos as métricas estruturais e a troca de vizinhos
(Tabela~\ref{tab:robustez}). \textbf{Em nenhuma variação a conclusão qualitativa muda}: a
clusterização é maior no bloco 1 que nos blocos 18 e 36, as distâncias encurtam a partir do
bloco 18 e a troca de vizinhos em relação à rede lexical cresce do bloco 1 ao 36.

\tabelaopcional{robustez}
    {Métricas-resumo em cada variação, ao lado da configuração principal. $\bar C$:
    clusterização local média; $\bar\ell$: distância média; $Q$: modularidade do Leiden;
    $J$: Jaccard médio com a rede lexical de mesmo $k$ e mesma regra de empate. Cada linha
    muda uma escolha em relação à principal.}
    {tab:robustez}

\paragraph{Número de vizinhos.} Com $k=5$ ou $k=20$, os valores mudam na direção esperada
(mais vizinhos, distâncias menores: no bloco 18, 7,3 com $k=5$, 5,1 com $k=10$ e 4,0 com
$k=20$), mas a ordem entre as camadas se mantém. A dominância relativa ao teto no bloco 36
fica entre 0,37 e 0,43 nos três valores de $k$.

\paragraph{Rede mútua.} Na versão em que a aresta exige escolha dos dois lados, a rede se
fragmenta e as distâncias crescem (17,0, 9,1 e 10,4 nos blocos 1, 18 e 36), mas a ordem entre as
camadas é a mesma.

\paragraph{Empates.} O desempate pela posição dá praticamente os mesmos valores que o sorteio
(Jaccard de 0,359 contra 0,354 entre \cod{lex} e \cod{L01}). Aceitar todos os empatados
(variante b) aumenta um pouco o Jaccard (0,41 no bloco 1), porque a vizinhança lexical fica
maior e contém mais cópias do próprio token, mas mantém a ordem. Com uma tolerância de empate
cem vezes maior ($\varepsilon=10^{-4}$), as métricas não mudam na segunda casa decimal, e as
vizinhanças de 2 a 4\% das linhas das redes contextuais mudam. A
Figura~\ref{fig:robustez-empates} separa as variantes por faixa de frequência. Duas delas
completam a leitura da P1: o Jaccard por composição de tipos ($J^w$), que ignora o sorteio,
fica perto de 1 entre \cod{lex} e \cod{L01} para os tipos frequentes, o que confirma que o
bloco 1 mantém os vizinhos do mesmo token; e a variante (c), que ignora o próprio tipo e compara
só quais \emph{outros} tipos estão por perto, mostra que esses outros tipos já mudam no bloco 1
(Jaccard entre 0,3 e 0,4 em todas as faixas).

\figuraopcional{robustez_empates}
    {Jaccard médio entre a rede lexical e cada camada nas regras de empate (sorteio, posição,
    todos os empatados e tipos distintos) e Jaccard por composição de tipos ($J^w$), por faixa
    de frequência.}
    {fig:robustez-empates}

\paragraph{Normalização final e centragem.} A saída normalizada do bloco 36 (\cod{L36n}),
que é o que o modelo usa para prever o próximo token, dá redes quase iguais às da saída bruta
(clusterização 0,343 contra 0,353; Jaccard com a lexical 0,094 contra 0,096). Subtrair dos
vetores a média de todas as ocorrências, o que remove a direção comum (a anisotropia), também
muda pouco (no bloco 36, distância média de 5,72 contra 5,28). A direção comum não é o que
produz os resultados.

\paragraph{Comunidades.} Com resoluções 0,5 e 2 do Leiden, e com o Louvain, a modularidade das
redes \cod{L18} e \cod{L36} fica entre 0,73 e 0,76, perto da principal (0,76 e 0,77).

\paragraph{O que não foi testado.} As concordâncias da P3 e a separação da P4 foram medidas só
na configuração principal (e nos 36 blocos), e não nas variações acima. Uma rede construída só
com os vértices do núcleo, prevista no plano, não foi feita: ela exige refazer o k-NN, porque o
k-NN de um subconjunto não é um recorte do k-NN completo (quando os vizinhos de fora saem, cada
vértice precisa escolher outros).
